% MPMC queue tradeoffs on asymmetric client silicon -- IEEEtran skeleton.
% Build: tectonic paper/main.tex   (figures come from scripts/make_plots.py)
%
% AUTHORSHIP NOTE (do not delete): drafts in this file are scaffolding produced
% with AI assistance. Before any submission the human author must verify every
% claim against paper/notes/claims.md and the data in paper/data/, rewrite the
% prose in their own voice, and check the venue's AI-disclosure policy.
\documentclass[conference]{IEEEtran}

\usepackage{graphicx}
\usepackage{booktabs}
\usepackage{amsmath}
\usepackage[hidelinks]{hyperref}

\newcommand{\figtag}{v1} % dataset tag; bump when re-running the matrix

% Placeholder box for figures that have not been generated yet.
\newcommand{\maybefig}[2]{%
  \IfFileExists{assets/#1_\figtag.pdf}%
    {\includegraphics[width=#2]{assets/#1_\figtag.pdf}}%
    {\fbox{\parbox[c][3.2cm][c]{0.9\linewidth}{\centering placeholder
      (\texttt{\detokenize{#1}}): run \texttt{scripts/reproduce.sh}}}}}

\begin{document}

\title{MPMC Queue Tradeoffs on Asymmetric Client Silicon:\\
A Characterization Study on Apple M2}

\author{\IEEEauthorblockN{Dedeepya Avancha}
\IEEEauthorblockA{Independent / [affiliation]\\
dedeepyaavancha@gmail.com}}

\maketitle

\begin{abstract}
Published evaluations of concurrent queue algorithms overwhelmingly report
results from symmetric x86 server processors with pinned threads and active
cooling. Client devices violate all three assumptions: cores are asymmetric
(performance and efficiency clusters), thread placement is delegated to the
OS (macOS exposes no affinity API on Apple silicon), and passive cooling makes
sustained load thermally sensitive. We characterize four classic
multi-producer/multi-consumer queue designs---a mutex-guarded deque, the
Michael--Scott lock-free queue with epoch-based reclamation, Vyukov's bounded
CAS queue, and a fetch-and-add ticket queue---plus an industrial reference
implementation, on an Apple M2 (4 performance + 4 efficiency cores). Our
experiments vary producer:consumer ratio, QoS-steered core placement, and
oversubscription, with a coordinated-omission-aware latency methodology and
thermally interleaved trials. [Headline findings inserted after the k=10 run:
(1) progress guarantees vs.\ tail latency under preemption; (2) how P/E
placement reshuffles throughput rankings; (3) where the textbook ranking
survives.] All results are reproducible with one command from our open
artifact.
\end{abstract}

\section{Introduction}
% Motivation skeleton -- rewrite in author's voice.
Lock-free queues are folklore-fast: the standard narrative says CAS beats
locks, FAA beats CAS, and progress guarantees are what you pay for with
complexity. That narrative was established on server-class symmetric
multiprocessors with pinned threads. Meanwhile, most code that ships runs on
client silicon: heterogeneous cores, opaque OS scheduling, passive cooling.

This paper asks a narrow, answerable question: \emph{how much of the textbook
MPMC-queue tradeoff story survives contact with an asymmetric, unpinnable,
thermally limited client SoC?} We make no claim of algorithmic novelty; the
contribution is a careful characterization and a reproducible methodology for
a measurement environment that prior work treats as out of scope.

Contributions:
\begin{itemize}
  \item A four-design comparison (blocking, lock-free node-based, lock-free
        array CAS, FAA ticket) implemented from the literature in a single
        harness, validated with ThreadSanitizer and invariant-checking stress
        tests, plus \texttt{moodycamel::ConcurrentQueue} as an industrial
        reference.
  \item A client-silicon methodology: QoS-class core steering (the only
        sanctioned placement control on macOS), oversubscription as a
        preemption proxy, coordinated-omission-aware paced latency, process-
        per-trial isolation, and thermal round-robin interleaving with a
        per-trial calibration probe on a fanless machine.
  \item [Findings, filled from data: H1 progress-guarantee tail flip under
        oversubscription; H2 P/E placement effects; the dedicated-core case
        where the mutex baseline is competitive or better.]
\end{itemize}

\section{Background and Related Work}
\label{sec:background}
% Condensed from paper/notes/related-work.md; keep citations tight.
\textbf{Designs.} Michael and Scott's two-lock and non-blocking queues
\cite{michael1996} anchor the node-based lineage; safe memory reclamation via
hazard pointers \cite{michael2004} or epochs \cite{fraser2004} is its
practical tax. Vyukov's bounded MPMC queue \cite{vyukov} popularized per-cell
sequence coordination. Fetch-and-add designs (LCRQ \cite{morrison2013})
exploit the observation that FAA scales better than contended CAS; the LMAX
Disruptor \cite{thompson2011} is the industrial ticket/sequencer ancestor.
Wait-free constructions \cite{yang2016} complete the progress taxonomy.

\textbf{Methodology.} We follow Hoefler and Belli's guidance on reporting
parallel-systems measurements \cite{hoefler2015} and Tene's critique of
latency measurement under coordinated omission \cite{tene}.

\textbf{Gap.} Prior evaluations assume symmetric cores, thread pinning, and
sustained-load thermal headroom; none of these hold on the devices where most
software runs. We are not aware of a systematic MPMC queue characterization on
asymmetric ARM client silicon.

\section{Queue Designs Under Study}
\label{sec:designs}
Table~\ref{tab:designs} summarizes the design space along the axes that the
experiments exercise.

\begin{table}[t]
\caption{Design space. ``Blocking-on-slot'': a preempted claimant blocks the
successor of its slot although claiming itself is wait-free.}
\label{tab:designs}
\centering
\small
\begin{tabular}{@{}lllll@{}}
\toprule
Design & Storage & Claim & Progress & Reclaim \\
\midrule
mutex+deque & -- & lock & blocking & none \\
Michael--Scott & nodes & CAS & lock-free & EBR \\
Vyukov & array & CAS & lock-free & none \\
FAA ticket & array & FAA & blk-on-slot & none \\
moodycamel & blocks & mixed & lock-free & internal \\
\bottomrule
\end{tabular}
\end{table}

% One paragraph per design: what it is, the one implementation subtlety that
% matters for interpretation (EBR stalls under preemption; Vyukov CAS retry;
% FAA irrevocable tickets), each ~4 sentences. Draft from learn/28.

\section{Methodology}
\label{sec:method}
\textbf{Platform.} Apple M2 (MacBook Air, passively cooled), 4 performance +
4 efficiency cores, macOS; Apple clang 17, \texttt{-O3}; 128-byte cache-line
padding on all hot atomics.

\textbf{Harness.} One process per trial emits one CSV row; the experiment
matrix interleaves queues round-robin within each trial round so cross-queue
comparisons occur under near-identical thermal state. A fixed
dependent-arithmetic calibration loop is timed before every trial as a
frequency proxy; its chronological trace (Fig.~\ref{fig:calib}) screens the
dataset for thermal drift. Trial~0 of each configuration is discarded as
warmup; we report medians with interquartile ranges over the remaining
trials \cite{hoefler2015}.

\textbf{Shutdown protocol.} A reserved poison-pill value terminates consumers
deterministically; this is required for the FAA design, whose irrevocable
tickets make counter-based termination deadlock-prone.

\textbf{Latency.} Producers are paced against a fixed schedule (total offered
load split across producers, staggered); the payload carries the
\emph{scheduled} send time, so producer lateness is charged to the reported
tail rather than silently omitted \cite{tene}.

\textbf{Placement and preemption.} macOS exposes no thread affinity on Apple
silicon; we steer placement with thread QoS classes
(\texttt{USER\_INTERACTIVE} biases performance cores,
\texttt{BACKGROUND} efficiency cores) and treat the result as a bias, not an
assignment. Oversubscription (thread count multiplied at fixed ratio) serves
as a controlled preemption generator.

\section{Results}
\label{sec:results}
% Every claim here must have a row in paper/notes/claims.md with data behind
% it. Numbers below are inserted from paper/data/summary_v1.md (k=5 triage)
% and finalized from the k=10 run.

\subsection{Throughput across producer:consumer ratios}
\begin{figure}[t]\centering
\maybefig{fig_throughput_ratio}{\linewidth}
\caption{Median throughput by P:C ratio, dedicated cores, default QoS.
Error bars: IQR.}
\label{fig:ratio}
\end{figure}
[Draft after data: where the mutex baseline sits; FAA vs CAS ordering; the
SPSC 1:1 reference point as the generality tax.]

\subsection{Oversubscription: progress guarantees under preemption (H1)}
\begin{figure}[t]\centering
\maybefig{fig_oversubscription}{0.9\linewidth}
\caption{Throughput at 4P:4C as threads are multiplied beyond cores
(log--log). }
\label{fig:oversub}
\end{figure}

\begin{figure}[t]\centering
\maybefig{fig_latency_tail}{\linewidth}
\caption{Paced-load latency percentiles, dedicated (left) vs.\ 4$\times$
oversubscribed (right), log scale.}
\label{fig:tail}
\end{figure}
[Draft after data: does the blocking-on-slot FAA design lose its dedicated-core
advantage under preemption, and by how much at p99.9?]

\subsection{Core placement via QoS (H2)}
\begin{figure}[t]\centering
\maybefig{fig_qos}{\linewidth}
\caption{Throughput at 4P:4C under QoS placement policies.}
\label{fig:qos}
\end{figure}
[Draft after data: E-core penalty per design; asymmetric producer/consumer
placements.]

\subsection{Thermal validity}
\begin{figure}[t]\centering
\maybefig{fig_calib}{\linewidth}
\caption{Calibration-probe trace across the full matrix (thermal screen).}
\label{fig:calib}
\end{figure}

\section{Threats to Validity}
\label{sec:threats}
Single machine and OS (stated scope: client characterization, not a general
ranking); QoS placement is a scheduler bias, not pinning -- we report policies,
not measured per-sample core residency; our implementations are faithful but
not micro-optimized to each design's best known form (the industrial reference
arm bounds this gap); Apple silicon's memory subsystem details are not
publicly documented, so microarchitectural attribution is necessarily
inferential.

\section{Conclusion}
[One paragraph, written last, matching only what the data supports.]

\section*{Artifact}
All code, data, and figure pipelines:
\url{https://github.com/Dd1235/SPSC_Queue}. \texttt{scripts/reproduce.sh}
rebuilds the dataset and every figure in this paper.

\bibliographystyle{IEEEtran}
\bibliography{refs}

\end{document}
